\section{Gobierno de datos maestros}
El gobierno de datos maestros (\textit{Master Data Management}, MDM) es la disciplina encargada de asegurar que las entidades centrales de una organización —clientes, proveedores, materiales, entre otras— se mantengan consistentes, correctas y libres de duplicidad a través de los distintos sistemas que las consumen. \citeA{dama2017} ubica el MDM como una de las áreas funcionales del gobierno de datos, y lo define como el conjunto de procesos que permiten mantener una versión autorizada y compartida de las entidades de negocio críticas. A diferencia de los datos transaccionales, que registran hechos puntuales, los datos maestros describen objetos persistentes que son referenciados por múltiples procesos a lo largo del tiempo; esa persistencia es precisamente lo que hace que un error introducido en el momento de la creación se propague y se vuelva costoso de revertir.

\subsection{Dimensiones de calidad del dato maestro}
La calidad de los datos no es un atributo único sino un constructo multidimensional. \citeA{wang1996} establecieron, a partir de un estudio empírico con consumidores de datos, un marco jerárquico que organiza la calidad en cuatro categorías —intrínseca, contextual, de representación y de accesibilidad— y demostraron que los usuarios manejan una concepción de calidad considerablemente más amplia que la de exactitud, que era la dimensión dominante en la literatura previa. Su distinción entre calidad intrínseca (el dato es correcto en sí mismo) y calidad contextual (el dato es adecuado para la tarea en que se usa) es especialmente pertinente en un catálogo de materiales, donde una descripción puede ser literalmente cierta y aun así resultar inservible para quien necesita distinguir dos repuestos similares.

Aplicando este marco al dato maestro específicamente, \citeA{loshin2009} señala que su calidad se mide típicamente en tres dimensiones operativas: \textit{unicidad} (ausencia de registros duplicados), \textit{exactitud} (que los atributos del registro describan correctamente el objeto real) y \textit{completitud} (que los campos obligatorios estén correctamente llenados). Estas tres dimensiones son las que estructuran el diagnóstico y la intervención del presente proyecto: el módulo de detección de duplicados atiende la unicidad, el módulo de clasificación atiende la exactitud, y el módulo de normalización atiende la completitud y la consistencia de representación.

\subsection{Impacto económico de la mala calidad de datos}
El deterioro de cualquiera de estas dimensiones tiene un efecto en cascada sobre los procesos que dependen del catálogo, como la planificación de compras, el control de inventario y el mantenimiento de activos. \citeA{redman1998} documentó que los efectos de la mala calidad de datos en una empresa típica no se limitan al costo directo de corregir los registros, sino que incluyen incremento del costo operativo, decisiones menos efectivas, deterioro de la capacidad de ejecutar la estrategia y erosión de la confianza organizacional en los sistemas de información. Esta distinción entre el costo visible de la corrección y el costo difuso de operar sobre datos defectuosos es la base conceptual sobre la que se construye la justificación financiera del presente trabajo: el ahorro atribuible a una herramienta de gobierno de datos proviene tanto del tiempo de gestión que evita como de los errores que impide introducir en primer lugar.

\subsection{Estándares para la descripción de datos maestros de materiales}
La descripción de un material en un catálogo industrial no es texto libre arbitrario: existen estándares internacionales que definen cómo debe estructurarse para que sea inequívoca y portable entre organizaciones y sistemas.

La norma \textit{ISO 8000} constituye la familia de estándares internacionales de calidad de datos, con varias partes dedicadas específicamente al dato maestro. En particular, la norma ISO 8000-110 especifica los requisitos para el intercambio de datos característicos de datos maestros en cuanto a sintaxis, codificación semántica y conformidad con una especificación de datos \cite{iso8000110}. El principio central es que un dato de calidad debe ser \textit{portable}: su significado debe estar codificado de forma explícita y no depender del sistema o de la convención interna de quien lo creó.

Complementariamente, la serie \textit{ISO 22745} define los diccionarios técnicos abiertos (\textit{Open Technical Dictionary}, OTD) y su aplicación a los datos maestros \cite{iso22745}. Bajo este enfoque, un material se describe mediante un identificador de concepto —el sustantivo o término principal que designa la clase de objeto— acompañado de un conjunto de pares propiedad-valor que lo especifican. Esta convención, conocida en la práctica de catalogación industrial como formato \textit{sustantivo-modificador} (\textit{noun-modifier}), es la que subyace al estándar corporativo empleado en la unidad de negocio de este proyecto, donde las descripciones siguen una estructura de término principal seguido de modificadores delimitados por separadores. La normalización automática que realiza GAMMA consiste, en esencia, en llevar una descripción propuesta en lenguaje libre por el gestor a esta forma estructurada.

Finalmente, para la clasificación de los materiales existe el \textit{United Nations Standard Products and Services Code} (UNSPSC), una taxonomía abierta y multisectorial de productos y servicios administrada por el Programa de las Naciones Unidas para el Desarrollo, organizada en una jerarquía de cuatro niveles —segmento, familia, clase y comodidad— codificada como un número de ocho dígitos \cite{undp2023}. El catálogo de la unidad de negocio mantiene una correspondencia entre sus clases internas de material y los códigos UNSPSC, lo que permite anclar la taxonomía propia a un referente externo estandarizado.

\section{Inteligencia de negocios y medición del desempeño del proceso}
Una intervención sobre un proceso operativo solo puede sostenerse en el tiempo si existe un mecanismo para medir su efecto. La inteligencia de negocios (\textit{Business Intelligence}, BI) aporta ese componente: el conjunto de arquitecturas, modelos y herramientas que transforman datos operativos en indicadores que soportan la decisión gerencial.

\subsection{Modelado dimensional}
\citeA{kimball2013} establecieron el modelado dimensional como el patrón de diseño predominante para estructurar datos orientados al análisis. Su propuesta organiza la información en tablas de hechos —que registran las mediciones numéricas de un proceso de negocio, con una granularidad explícitamente declarada— y tablas de dimensión —que contienen los atributos descriptivos por los que esas mediciones se filtran y agrupan. La ventaja de este esquema, conocido como \textit{estrella}, frente a un modelo normalizado, es que privilegia la comprensibilidad para el usuario de negocio y el desempeño de consulta por encima de la ausencia de redundancia. En el presente proyecto, este patrón orienta el diseño de la capa de consumo, donde cada solicitud de creación de material constituye el grano de la medición y los atributos de gestor, tipo de material y período conforman las dimensiones de análisis.

\subsection{Indicadores clave de desempeño}
Un indicador clave de desempeño (\textit{Key Performance Indicator}, KPI) es una medida cuantificable, asociada a un objetivo declarado, que permite evaluar si un proceso avanza en la dirección esperada. \citeA{kimball2013} advierten que la utilidad de un indicador depende críticamente de que su definición —numerador, denominador, granularidad temporal y población incluida— esté documentada de forma no ambigua, ya que una misma métrica calculada bajo criterios distintos produce cifras irreconciliables entre áreas. Esta consideración es la que motiva, en este trabajo, la definición explícita de la población de medición y la exclusión sistemática del tráfico de pruebas del cálculo de los indicadores.

\subsection{Valoración económica de la información}
\citeA{laney2018} propone bajo el término \textit{infonomics} la tesis de que la información debe tratarse como un activo económico formal, sujeto a valoración, gestión y monetización con el mismo rigor que se aplica a los activos físicos. El autor distingue entre la monetización directa —vender o intercambiar información— y la indirecta, que consiste en aprovechar la información para reducir costos, mitigar riesgos o mejorar la eficiencia de un proceso interno. Es esta segunda modalidad la que corresponde al presente proyecto: el valor generado no proviene de comercializar el catálogo de materiales, sino de que un catálogo de mayor calidad reduce el tiempo de gestión y evita los costos aguas abajo asociados a registros duplicados o mal clasificados.

\section{Metodología para proyectos de analítica}
El desarrollo de un proyecto de analítica requiere un marco de proceso que ordene sus fases y haga explícita la naturaleza iterativa del trabajo. \citeA{wirth2000} formalizaron el modelo \textit{CRISP-DM} (\textit{Cross-Industry Standard Process for Data Mining}), concebido como un proceso de referencia independiente tanto del sector industrial como de la tecnología empleada. El modelo organiza el trabajo en seis fases —comprensión del negocio, comprensión de los datos, preparación de los datos, modelado, evaluación y despliegue— y establece que la secuencia entre ellas no es estrictamente lineal: los hallazgos de una fase posterior frecuentemente obligan a revisar decisiones tomadas en una anterior. CRISP-DM continúa siendo el modelo de proceso de referencia más extendido para proyectos de minería de datos y ciencia de datos aplicada, y es el que estructura la metodología seguida en este trabajo.

\section{Arquitectura de datos por capas (medallón)}
La arquitectura medallón es un patrón de diseño para \textit{pipelines} de datos que organiza el flujo de información en capas sucesivas de refinamiento, denominadas Bronce, Plata y Oro \cite{databricks2026}. La capa Bronce conserva el dato tal como llega de la fuente original, sin transformaciones. La capa Plata aplica limpieza, normalización y validaciones de calidad. La capa Oro contiene datos ya agregados o modelados específicamente para su consumo final, ya sea por un tablero de indicadores, un modelo analítico o una aplicación. Este patrón permite reprocesar el pipeline completo ante un error sin perder el dato original, y desacopla la lógica de transformación de la lógica de consumo.

Este esquema por capas se inscribe en la propuesta arquitectónica del \textit{lakehouse}, formulada por \citeA{zaharia2021}, que busca unificar en una sola plataforma las capacidades de gestión y confiabilidad propias del almacén de datos con la flexibilidad y el soporte a cargas de ciencia de datos característicos del lago de datos. Los autores argumentan que la separación tradicional entre ambos entornos genera problemas de obsolescencia del dato, duplicación de almacenamiento y costo total de propiedad, y que una arquitectura de capas sobre formatos abiertos con garantías transaccionales resuelve simultáneamente los requerimientos de reportería y de modelado analítico. En el presente proyecto este patrón se implementa sobre un motor relacional, aprovechando que el volumen del catálogo permite prescindir de una plataforma distribuida sin sacrificar la separación lógica entre capas.

\section{Procesamiento de lenguaje natural para normalización de texto}
Las descripciones de materiales en un catálogo industrial constituyen un tipo particular de texto: breve, técnico, con alta densidad de abreviaturas y sin una gramática formal. Se trata además de texto sujeto a una restricción estructural fuerte, ya que los sistemas ERP imponen un límite estricto de caracteres al campo de descripción corta, lo que induce a los gestores a abreviar de formas idiosincráticas y no documentadas.

\subsection{Enfoques basados en reglas y basados en modelos de lenguaje}
El procesamiento de lenguaje natural (\textit{Natural Language Processing}, NLP) ofrece dos familias de enfoques para trabajar con este tipo de texto: una basada en reglas y diccionarios de abreviatura, precisa pero de mantenimiento manual constante y de cobertura limitada al vocabulario previamente registrado; y otra basada en modelos de lenguaje de gran escala (\textit{Large Language Models}, LLM), que aprovecha el conocimiento general del lenguaje adquirido durante el entrenamiento para inferir la forma estándar de un término técnico sin necesidad de un diccionario exhaustivo.

\subsection{La arquitectura Transformer}
Los LLM actuales se construyen sobre la arquitectura \textit{Transformer}, introducida por \citeA{vaswani2017}, que reemplazó los mecanismos recurrentes por un mecanismo de atención capaz de ponderar la relevancia de cada palabra de una secuencia respecto a las demás. Esta arquitectura permite procesar la secuencia completa en paralelo y capturar dependencias entre términos distantes, lo que resultó determinante para el escalamiento de los modelos de lenguaje a los tamaños actuales.

\subsection{Modelos de lenguaje aplicados a tareas de preparación de datos}
\citeA{brown2020} demostraron que, a partir de cierta escala, los modelos de lenguaje adquieren la capacidad de resolver tareas nuevas a partir de unas pocas demostraciones incluidas en la propia instrucción (\textit{few-shot learning}), sin necesidad de reentrenamiento ni de ajuste de parámetros. Esta propiedad es la que hace viable aplicar un LLM a la normalización de descripciones sin construir un conjunto de entrenamiento etiquetado específico para la tarea.

La pertinencia de este enfoque para problemas de datos ha sido evaluada de forma sistemática. \citeA{narayan2022} formularon cinco tareas clásicas de limpieza e integración de datos —entre ellas la imputación de valores faltantes, la corrección de errores y el emparejamiento de entidades— como tareas de \textit{prompting} sobre modelos fundacionales, y encontraron que estos alcanzan un desempeño competitivo con el estado del arte a pesar de no haber sido entrenados específicamente para dichas tareas. Este resultado respalda empíricamente la decisión de delegar la normalización de descripciones a un modelo de lenguaje en lugar de a un motor de reglas, que es el enfoque adoptado en el módulo correspondiente de GAMMA.

\section{Detección de duplicados y resolución de entidades}

\subsection{El problema de la detección de duplicados}
La existencia de múltiples representaciones de una misma entidad dentro de una base de datos es un problema clásico y ampliamente estudiado. \citeA{elmagarmid2007} lo caracterizan como la situación en que registros duplicados carecen de una clave común y presentan además errores de transcripción, información incompleta o ausencia de formatos estándar, lo que impide resolverlos mediante comparación exacta. Los autores sistematizan las familias de métricas de similitud aplicables —basadas en caracteres, en tokens o en modelos probabilísticos— y las técnicas para hacer escalable la comparación, cuyo costo es cuadrático en el número de registros si se realiza de forma exhaustiva. Su caracterización aplica de forma directa al catálogo de materiales estudiado, donde las descripciones libres, las abreviaturas inconsistentes entre gestores y el límite de caracteres del campo producen exactamente el tipo de variación que hace fallar la comparación literal.

\subsection{Similitud por q-gramas y trigramas}
La búsqueda por similitud difusa (\textit{fuzzy matching}) resuelve este problema descomponiendo cada cadena de texto en subsecuencias de caracteres consecutivos, denominadas \textit{q-gramas}, y comparando el grado de superposición entre los conjuntos de q-gramas de dos cadenas. Cuando la longitud de la subsecuencia es tres, se les denomina \textit{trigramas} \cite{postgresql2026}. Este enfoque es tolerante a variaciones menores de escritura, a la inversión del orden de las palabras y al uso de separadores distintos, y no requiere que las cadenas comparadas compartan una estructura gramatical común.

\citeA{gravano2001} demostraron que este tipo de emparejamiento aproximado puede expresarse mediante operaciones relacionales convencionales y ejecutarse dentro del propio motor de base de datos, aprovechando sus índices y su optimizador sin necesidad de infraestructura externa. La extensión \texttt{pg\_trgm} de PostgreSQL \cite{postgresql2026} implementa este principio de forma nativa, permitiendo indexar el catálogo completo y ejecutar búsquedas de similitud con baja latencia. Esta propiedad es determinante para el diseño de GAMMA, ya que la detección de duplicados debe resolverse de forma interactiva, en el momento en que el gestor propone una descripción.

\subsection{Umbral de similitud y compromiso entre exhaustividad y precisión}
Un motor de similitud difusa no produce una decisión binaria sino un puntaje continuo, por lo que su operación requiere fijar un umbral a partir del cual dos registros se consideran candidatos a duplicado. La elección de ese umbral define un compromiso explícito: un valor bajo maximiza la exhaustividad —detecta más duplicados reales— a costa de generar falsos positivos que el gestor debe descartar manualmente; un valor alto reduce las alertas espurias pero deja pasar duplicados con redacción divergente \cite{elmagarmid2007}. Dado que el umbral determina simultáneamente la utilidad operativa de la herramienta y la magnitud del problema que se reporta al medir el catálogo, su valor debe declararse explícitamente junto con cualquier cifra de duplicidad, y su efecto debe caracterizarse mediante un análisis de sensibilidad.

\section{Representación vectorial de texto}
Antes de aplicar un algoritmo de aprendizaje automático sobre texto, este debe convertirse en una representación numérica.

\subsection{TF-IDF}
Una de las técnicas más establecidas es \textit{TF-IDF} (\textit{Term Frequency – Inverse Document Frequency}), propuesta por \citeA{salton1988}, que pondera cada término según su frecuencia relativa dentro de un documento y su rareza a través de todo el corpus. La intuición es que un término frecuente en un documento pero infrecuente en la colección resulta altamente discriminante, mientras que uno presente en casi todos los documentos aporta poca información para distinguirlos. Esta representación es la que se emplea como entrada de los algoritmos de clasificación evaluados en este proyecto.

\subsection{N-gramas de carácter}
La variante empleada en este trabajo no opera sobre palabras completas sino sobre subcadenas de caracteres. \citeA{cavnar1994} demostraron que la categorización de texto basada en frecuencias de n-gramas de carácter es robusta frente a errores tipográficos y variaciones morfológicas, precisamente porque un error puntual afecta únicamente a los pocos n-gramas que lo contienen y deja intacto el resto de la representación. Esta propiedad es especialmente valiosa en un catálogo de materiales, donde las descripciones contienen abreviaturas no normalizadas, unidades de medida, fracciones y errores de digitación que fragmentarían una representación basada en palabras completas, pero que degradan solo parcialmente una representación basada en caracteres.

\section{Algoritmos de clasificación supervisada evaluados}
La asignación automática de una categoría a partir de la descripción de un material es un problema de clasificación de texto supervisada. \citeA{sebastiani2002} consolidó el enfoque inductivo —construir el clasificador a partir de ejemplos previamente etiquetados— como el paradigma dominante en esta tarea, frente a la definición manual de reglas por expertos de dominio. Dada la diversidad de familias de algoritmos disponibles, se evaluaron seis enfoques distintos sobre el mismo conjunto de datos, descritos a continuación.

\subsection{Regresión logística}
Modelo lineal que estima la probabilidad de pertenencia a una clase mediante una función logística aplicada a una combinación lineal de las variables de entrada. Fue formalizado por \citeA{cox1958} como un método de regresión para variables de respuesta binaria, y se ha extendido posteriormente a problemas multiclase.

\subsection{Máquinas de vectores de soporte (SVM lineal)}
Introducidas por \citeA{cortes1995}, buscan el hiperplano que separa las clases maximizando el margen entre los puntos más cercanos de cada una. En su variante lineal (\texttt{LinearSVC}), el modelo opera directamente sobre la representación vectorial del texto sin necesidad de una función núcleo (\textit{kernel}), lo que la hace especialmente eficiente en espacios de alta dimensionalidad como los generados por TF-IDF.

\subsection{Random Forest}
Propuesto por \citeA{breiman2001}, es un método de ensamble que combina múltiples árboles de decisión entrenados sobre subconjuntos aleatorios de datos y variables, y agrega sus predicciones por votación. Su capacidad de capturar relaciones no lineales lo hace robusto frente a ruido en los datos, a costa de un mayor costo computacional que un modelo lineal.

\subsection{XGBoost}
\citeA{chen2016} presentaron este sistema de \textit{gradient boosting} sobre árboles de decisión, que construye los árboles de forma secuencial, corrigiendo en cada iteración los errores del árbol anterior. Es reconocido por su desempeño competitivo en tareas de clasificación estructurada, aunque su tiempo de entrenamiento crece considerablemente con la dimensionalidad del vocabulario.

\subsection{fastText}
Desarrollado por \citeA{joulin2017}, es un clasificador lineal que representa cada documento como el promedio de los vectores de sus palabras (o subpalabras), entrenado de forma conjunta con la tarea de clasificación. Fue diseñado explícitamente para ofrecer un desempeño comparable al de modelos de aprendizaje profundo con tiempos de entrenamiento varios órdenes de magnitud menores.

\subsection{Transformer (embeddings contextuales)}
Los modelos basados en la arquitectura \textit{Transformer} \cite{vaswani2017} generan representaciones vectoriales del texto que capturan el significado de cada palabra en función de su contexto, a diferencia de TF-IDF, que trata cada término de forma independiente. Estas representaciones, preentrenadas sobre grandes volúmenes de texto general, suelen requerir un proceso adicional de ajuste fino (\textit{fine-tuning}) para adaptarse a vocabularios técnicos especializados como el de un catálogo industrial. Este requerimiento implica que su desempeño depende de forma crítica del presupuesto de cómputo y de la configuración del entrenamiento, un factor que debe considerarse al interpretar los resultados de una comparación entre familias de algoritmos.

\section{Evaluación de clasificadores multiclase}
La comparación entre algoritmos exige definir con precisión las métricas empleadas, particularmente cuando el problema involucra un número elevado de clases con distribución desigual.

\subsection{Métricas y esquemas de promediado}
\citeA{sokolova2009} realizaron un análisis sistemático de las medidas de desempeño utilizadas en tareas de clasificación binaria, multiclase, multietiqueta y jerárquica, caracterizando qué cambios en la matriz de confusión deja invariante cada medida. Su análisis hace explícito que la elección de la métrica no es neutral: la \textit{exactitud} (\textit{accuracy}) global queda dominada por las clases mayoritarias, mientras que las medidas promediadas ofrecen lecturas distintas del mismo resultado. El promedio \textit{macro} calcula la métrica de forma independiente para cada clase y luego promedia sin ponderar, otorgando el mismo peso a una clase con miles de ejemplos que a una con tres; el promedio \textit{ponderado} (\textit{weighted}) pondera cada clase por su frecuencia, aproximándose al comportamiento de la exactitud global. La brecha entre ambos promedios constituye por sí misma un diagnóstico del desbalance del problema, por lo que en este trabajo se reportan ambos.

\subsection{Exactitud top-k}
En un escenario de asistencia al usuario, donde el sistema propone sugerencias que una persona valida, la métrica relevante no es únicamente si la primera predicción es correcta, sino si la respuesta correcta aparece entre las primeras $k$ opciones presentadas. La exactitud \textit{top-k} mide esta propiedad y resulta más representativa del valor operativo del sistema cuando la interfaz muestra varias alternativas al gestor, ya que refleja la proporción de casos en que la herramienta le evita realizar la búsqueda manual.

\subsection{Desbalance de clases}
\citeA{he2009} revisaron de forma exhaustiva la problemática del aprendizaje sobre datos desbalanceados, donde la distribución desigual de ejemplos entre clases sesga a los algoritmos hacia las clases mayoritarias y degrada el desempeño sobre las minoritarias, que con frecuencia son las de mayor interés. Los autores sistematizan las estrategias disponibles —remuestreo, aprendizaje sensible al costo y métodos de ensamble— y advierten sobre la insuficiencia de la exactitud global como criterio de evaluación en estos escenarios. Esta consideración es central en el presente proyecto, dado que el catálogo estudiado presenta una distribución fuertemente asimétrica, con un número considerable de clases representadas por muy pocos materiales.

\subsection{Protocolo de validación}
Para obtener una estimación no sesgada del desempeño, el conjunto de datos se particiona en subconjuntos de entrenamiento y prueba, de modo que la evaluación se realice exclusivamente sobre ejemplos no vistos durante el ajuste del modelo. En problemas con clases desbalanceadas se emplea una partición \textit{estratificada}, que preserva la proporción de cada clase en ambos subconjuntos y evita que las clases minoritarias queden ausentes del conjunto de prueba \cite{he2009}. Cuando el modelo seleccionado se reentrena posteriormente sobre la totalidad de los datos para su despliegue, la métrica reportada continúa siendo la obtenida sobre la partición de prueba, y así debe declararse para no atribuir al modelo en producción una estimación medida sobre una configuración distinta.

\section{Calibración de probabilidades y predicción selectiva}
Un sistema de asistencia que sugiere una categoría necesita acompañar cada sugerencia de una medida de confianza interpretable, tanto para decidir cuándo presentarla como para permitir al usuario ponderar cuánto fiarse de ella.

\subsection{Calibración de probabilidades}
Un clasificador está \textit{calibrado} cuando las probabilidades que emite corresponden a frecuencias empíricas reales: de los casos a los que asigna una confianza de 0.80, aproximadamente el 80\% deberían resultar correctos. Las máquinas de vectores de soporte no producen probabilidades de forma nativa, sino distancias al hiperplano separador. \citeA{platt1999} propuso resolver esta limitación ajustando una función sigmoide sobre las salidas del modelo, de manera que las distancias se transformen en estimaciones de probabilidad conservando la propiedad de dispersión del SVM original. \citeA{zadrozny2002} extendieron posteriormente el tratamiento a problemas multiclase, mostrando que estos pueden descomponerse en múltiples problemas binarios cuyas estimaciones se combinan. Esta técnica es la que permite emplear un clasificador de margen máximo dentro de un flujo que requiere probabilidades, como es el caso de la arquitectura adoptada en este proyecto.

\subsection{Clasificación con opción de rechazo}
Disponer de probabilidades calibradas habilita una estrategia de \textit{predicción selectiva}: en lugar de emitir siempre una predicción, el sistema puede abstenerse cuando su confianza es insuficiente y derivar el caso al criterio humano. \citeA{chow1970} formalizó este planteamiento derivando la relación óptima entre la tasa de error y la tasa de rechazo para un sistema de reconocimiento bayesiano, y estableciendo que existe un compromiso cuantificable entre ambas: elevar el umbral de confianza reduce el error sobre los casos aceptados a costa de reducir la cobertura del sistema.

Este marco es el que fundamenta la operación de GAMMA como herramienta de asistencia y no de automatización plena. Al caracterizar empíricamente la curva de exactitud frente a cobertura para distintos umbrales de confianza, es posible seleccionar un punto de operación que refleje la tolerancia al error de la organización, y sustentar con datos la decisión de que el gestor conserve el control de la creación final del material.
